% Independent, scoped audit. No existing source is changed.
\section*{Lemma 5.1 with a polynomially large residue: a scoped derivation}
\label{sec:l51-large-a-audit-20261009}

This note audits the actual use of Lemma 5.1 in the retained-phase
proposal.  It does not certify Proposition 5.2, its preliminary dispersion
lemmas, or a prime-counting theorem.  The spectral theorems identified below
remain external inputs.  All occurrences of a power of $a$ mean $|a|$.
The source is the current local text of Lichtman's 2023 paper, lines
1911--2410; references below are to that text unless otherwise specified.

\paragraph{The precise restricted conclusion.}
Let $0<\sigma$, $0\leq\kappa\leq\sigma/100$, and $NM\asymp x$.
Use exactly the coefficients, cross-coprime set, product endpoints and fixed
$d$ of Lemma 5.1, with $(d,a)=1$, but separate the residue range
$0<|a|\leq x^{1+\kappa}$ from the source margin.  Put
\[
 D_0={N^2R^2S\over x},\qquad
 D_1={|a|^\theta NR^2S^5Q\over x^2},\qquad
 D_2={(NR/S)^{2\theta}NR^2S^5Q\over x^2},\quad \theta=7/32.
\]
Assume $Q<Nx^{-\sigma}$, $D_0\leq x^{-7\sigma}$ and
$D_1,D_2\leq x^{-14\sigma}$.  The proof gives the useful natural-scale
estimate
\begin{equation}\label{eq:l51-natural-large-a}
 |\mathcal W|\ll {N^2\over Q}x^{o(1)}
 \left\{D_0^{1/2}+\big((1+|a|/x)D_1+D_2\big)^{1/4}\right\}.
\end{equation}
Here the $o(1)$ also denotes independently selectable losses from the
spectral and smoothing parameters, with all scales polynomially bounded
in $x$.  More precisely, their total contribution to the exponent on the
right can be made at most $\sigma/4$.  After the additional retained-phase
variation loss $x^{5\kappa+o(1)}$, the saving is at least
\[
 {7\sigma\over2}-{21\kappa\over4}-{\sigma\over4}>3\sigma.
\]
Consequently the bound $N^2/(Qx^{\sigma})$ needed by the proposed bridge
has spare margin.  This conclusion uses the coprime spectral subclass
proved in the next two paragraphs; it does not use Lemma 5.1's conclusion
as an input.

\paragraph{What the cited Fourier-coefficient result actually says.}
The literal same-$j$ inequality printed as Lemma 9.6 (lines 4808--4830)
is not a literal restatement of Assing--Blomer--Li, Lemma 3.2.
In the primary text of \emph{Uniform Titchmarsh divisor problems},
arXiv:2005.13915v1, pp. 8--11, the proof requires $(a,rs)=1$, treats an
orthonormal oldform space, and retains a cusp phase.  In the $q=K=1$
specialization its modified coefficient is
\begin{equation}\label{eq:l51-abl-phase}
 b^*_{n,r,s}=\mathbf1_{(n,a)=1}b'_{nn'',r,s}
 e\left(\pm {n(an''-1)\overline s\over r}\right),\qquad n''\mid a^\infty.
\end{equation}
Thus $\|b^*(n'')\|_2\leq\|\widetilde b(n'')\|_2$.
In fact the source's earlier coefficient
$b'_n=e(\mp an\overline s/r)b_n$ cancels the $an''$ part, leaving a
unit-modulus phase $e(\mp n\overline s/r)$.  Cancellation is convenient,
but the norm inequality alone suffices.  No phase is deleted by an
approximation and no small-$a$ hypothesis is used here.

For clarity, the needed substitute for the printed same-$j$ assertion
is a Hilbert-space assertion.  Fix a level $q=rs$ and one newform $f$;
let $E_f$ be the space spanned by its oldforms at level $q$.  Normalize
Fourier coefficients as in the Kuznetsov pairing (equivalently absorb
$\sqrt{|n|}$ in the Assing--Blomer--Li convention).  The orthonormal basis
expansion (3.3) in that paper has divisor-bounded coefficients.  Apply its
Lemma 3.1 to each normalized oldform $f(dz)$.  Because $(a,q)=1$, the
integer $K$ in that lemma is $1$.  The Hecke factor is bounded by
$(|a|n'')^{\theta/2+\delta}$, for any fixed $\delta>0$; its remaining
phase is independent of the oldform label $d$.  For each normalized
$f(dz)$, complete it to an orthonormal basis of $E_f$.  The square of any
one resulting linear functional is at most its full squared norm.
Summing the divisor-bounded number of labels and taking a square root
therefore gives
\begin{equation}\label{eq:l51-packet}
 \|B_a\|_{E_f}\ll_\delta (|a|qN)^\delta
 \sum_{n''\mid a^\infty}(|a|n'')^{\theta/2}
 \|B^*(n'')\|_{E_f}.
\end{equation}
The same argument applies to the continuous spectrum; holomorphic
forms have at least this bound.  This is the oldform argument of the
cited source, applied within a single packet, rather than a new assertion
of multiplicativity for an arbitrary basis vector.

\paragraph{Why the interval side and the exceptional weights survive.}
In the Kuznetsov bilinear pairing write $A_f$ for the vector coming from
$\sum_{m\sim M}a_m\rho_{j\infty}(m)$ and $B_{a,f}$ for the other vector.
Apply Cauchy--Schwarz to the pairing within $E_f$, then
\eqref{eq:l51-packet}, and only then sum over $f$ and levels.  Every
oldform in $E_f$ has the same spectral parameter.  Hence an arbitrary
nonnegative weight depending on that parameter commutes with this
operation.  For each $n''$ one can choose
$X_0=1+RSn''/N$ separately and apply Cauchy--Schwarz with the two weights
used in (10.8)--(10.9).  An absolute constant covers the source's split
of $1+X^{2i\kappa_j}$.  No coefficient depending on $m$ has been inserted:
the squared norm of $A_f$ is invariant under the orthonormal changes of
basis above.  In particular, the interval-coefficient hypothesis of
Theorem 9.5 remains available on the $m$ side.

On the $n$ side Theorem 9.1 permits arbitrary complex coefficients,
including \eqref{eq:l51-abl-phase}, separately at each level.
With $N'=N/n''$ and $\mu(1/s)=1/(rs)$, its first factor is governed by
$\mu N'X_0\asymp1+N'/(RS)$, and its second factor by
$N'/(RS)$.  It therefore gives
$x^{o(1)}(1+N'/(RS))\|\widetilde b(n'')\|_2^2$.
This also repairs the explanatory line 5157, which drops the factor
$\mu N'$ before replacing $X_0$ by its reciprocal-scale expression.
Theorem 8.2 deals with the regular spectrum.  Theorems 9.1 and 9.5 are
$a$-free spectral inputs proved in the cited paper; this note checks
their use, not their complete proofs.  Theorem 8.2 and Kuznetsov are
the Deshouillers--Iwaniec inputs cited there.  The stated spectral results
are the inputs used here.

The standard Poisson completion from Theorem 10.2 to Theorem 1.8 is used
only for $q_0=1$, with the coprime coefficients described next.  This
audit does not claim the printed version of Theorem 1.8 for unrestricted
coefficients with $(a,rs)>1$.  In particular it does not silently promote
the weaker polynomial in the earlier Assing--Blomer--Li theorem to the
sharper exceptional-spectrum polynomial of Lichtman.

\paragraph{The actual coprimality, fixed $d$, and counting input.}
The two $s$ variables produced by expanding the square in Lemma 5.1
come from an inner sum with $(s,a)=1$.  Retain this restriction even
when other positive restrictions are dropped.  Then the spectral
coefficient is supported on $z=s_1s_2$ with $(z,a)=1$.
The substitution into Theorem 1.8 is
$(C,D,N,R,S,q_0)=(C,B,Y,Z,1,1)$, so precisely the coprime subclass
$(a,rs)=1$ above is used.  Multiplying by arbitrary unit phases to remove
the absolute value of the inner exponential sum changes no norm bound.

The original statement already permits a fixed $d$: it has
$n_1\equiv n_2\pmod{qd}$ and the inverse of $dq$ in its phase
(lines 1935--1949).  Use $fdq=n_1-n_2$, so $|f|\leq2N/(dQ)$.
The congruence modulo $df$ can be discarded only after the positive
Cauchy--Schwarz square is formed.  Keeping this proof avoids any
unsupported claim that Lemma 5.1 accepts arbitrary weights in $q$.
In particular no change $q'=qd$ followed by an unproved sparse-modulus
extension is needed.

The remaining multiplicity estimate is
\[
 \|\widetilde b(n'')\|_2^2\ll x^{o(1)}(HS)^2F/n'',\quad F\ll N/Q.
\]
It is uniform in $a$, which occurs only in $n''\mid a^\infty$.
One may replace the source's unnecessary claims $k\asymp HS$ by
$0<|k|\ll HS$.  For the equal-$h$ contribution, the number of nonzero
multiples of $n''/c$ in that range is $O(HSc/n'')$, and the other choices
cost $HSx^{o(1)}$.  The factor $F/c$ gives $(HS)^2F/n''$.
For unequal $h$, write $d_1=(s_1,s_2)$, reduce the two $s$ variables,
and let $g$ be the gcd of the two nonzero $h$ differences.  There are
$O(Hd_1/S)$ possible $g$, $O(HS/d_1)$ choices of one $h,s$ pair,
and a divisor-bounded number for the other pair once the nonzero $k$
is specified.  Counting its required multiples gives
\[
 x^{o(1)}{H^3SF\over n''}
 \sum_{d_1\ll S}{(d_1,n'')\over d_1}
 \ll x^{o(1)}H^3SF/n''.
\]
Here $H/S\leq x^{o(1)}(Q/N)D_0<1$, so this is no larger than the
equal-$h$ contribution.  This argument needs no lower dyadic bound
on $k$ and no $a\ll x^\varepsilon$ restriction.

\paragraph{Scale calculation, retaining every large-$a$ term.}
After the outer Cauchy--Schwarz step (line 2044) the factor is
$x^{o(1)}N/(QRS^{3/2})$.  Take upper scales
\[
 B=NR,\quad C=NRS,\quad Z=S^2,\quad
 Y=NR^2S^3/M,\quad F=N/Q,\quad H=x^{o(1)}QR^2S^2/M.
\]
The diagonal square is at most $x^{o(1)}HSBFC$.
For the off-diagonal, the spectral polynomial and the multiplicity
bound give
\[
 J^2\ll x^{o(1)}\left[
 |a|^\theta(CBZ+|a|Y)Z+
 C^{1+2\theta}BZ^{2-2\theta}+B^2YZ\right](HS)^2F.
\]
Indeed $Y/Z=D_0$ and $B^2<x$; and the correct identity is
$CBZ\asymp xY=N^2R^2S^3$, not the $S^5$ expression printed at line
2377.  Thus $B^2Y$ is absorbed into $CBZ$, and retaining $|a|Y$
produces the factor $1+|a|/x$ in
\eqref{eq:l51-natural-large-a}.  The off-diagonal contributes through
two square roots, which is why this extra factor costs $\kappa/4$,
not $\kappa$ in the final estimate.

Uniformly for $|a|\leq x^A$ and fixed $\theta<1$,
$\sum_{n''\mid a^\infty}(n'')^{\theta-1}=x^{o(1)}$.
For example its Euler product is bounded by splitting primes at
$\log x$: the logarithm is $O((\log x)^\theta)=o(\log x)$.
Truncated counts of such $n''$ are also $x^{o(1)}$, by Rankin's trick
with an arbitrarily small fixed exponent.  Therefore neither divisor
factor hides a condition $a\ll x^\varepsilon$.

\paragraph{What is, and is not, closed.}
The needed coprime large-residue use of Lemma 5.1 has no surviving
small-$a$ obstruction in the tracked proof after the packet-level
correction above.  Its natural power-saving bound and a separate
$\sigma,\kappa$ ledger are explicit.  This is a derivation relative to
the named spectral inputs and standard completion, not an independent
reproof of those theorems.  The full Proposition 5.2 repair, including
its preliminary errors, fixed-$d$ summation, and final normalization,
is outside this report.  The project-wide prime conclusions remain
unresolved until those interfaces and any separate weighted-inequality
issue have been audited.
